\par\noindent\textbf{(1)} 轴对称Laplace解可写 \(\Phi_{\rm in}=\sum_l A_lr^lP_l(\cos\theta)\)、\(\Phi_{\rm out}=\sum_l B_lr^{-l-1}P_l(\cos\theta)\)。只需 \(l=0,2\) 两阶：其他模式边界源为零，正交性保证对应系数均为零。
\par\noindent\textbf{(2)} 球面电势连续给 \(A_lR^l=B_lR^{-l-1}\)，法向Gauss跳跃 \(-\varepsilon_0(\partial_r\Phi_{\rm out}-\partial_r\Phi_{\rm in})|_R=\sigma_lP_l\)。代入得到 \(\varepsilon_0(2l+1)A_lR^{l-1}=\sigma_l\)，于是 \(A_l=\sigma_lR^{1-l}/[\varepsilon_0(2l+1)]\)、\(B_l=\sigma_lR^{l+2}/[\varepsilon_0(2l+1)]\)。
\par\noindent\textbf{(3)} 因此 \(\Phi_{\rm in}=\sigma_0R/\varepsilon_0+\sigma_2r^2P_2(\cos\theta)/(5\varepsilon_0R)\)，\(\Phi_{\rm out}=\sigma_0R^2/(\varepsilon_0r)+\sigma_2R^4P_2(\cos\theta)/(5\varepsilon_0r^3)\)。轴上 \(P_2(1)=1\) 且 \(E_z=-\partial_z\Phi_{\rm in}=-2\sigma_2z/(5\varepsilon_0R)\)，四极场内部呈线性。
\par\noindent\textbf{(4)} 总自由电荷 \(Q=\int\sigma\,dA=4\pi R^2\sigma_0\)，因为 \(\int P_2d\Omega=0\)。若 \(\sigma_2=0\)，内部场为零，外部 \(\Phi=Q/(4\pi\varepsilon_0r)\)；若 \(\sigma_0=0\)，远场从 \(r^{-3}\) 四极项开始而非单极库仑项。
